\begin{table}[t]
\caption{Total Effect Without Topic Adjustment — H1 (Probabilistic Time)} 
\fontsize{12.0pt}{14.0pt}\selectfont
\begin{tabular*}{\linewidth}{@{\extracolsep{\fill}}lcc}
\toprule
  & (1) Main‡ & (2) Total‡ \\ 
\midrule\addlinespace[2.5pt]
Invasion & -0.217*** & -0.240*** \\ 
 & (0.036) & (0.035) \\ 
Topic FEs & X & - \\ 
Source Z-Score & X & X \\ 
N (within bandwidth) & 45,366 & 45,398 \\ 
Clusters & — & — \\ 
\bottomrule
\end{tabular*}
\begin{minipage}{\linewidth}
\vspace{.05em}
+ p < 0.1, * p < 0.05, ** p < 0.01, *** p < 0.001\\
The total-effect column removes the topic adjustment from the corresponding
specification and is a distinct estimand, not a robustness check: it includes
shifts in the topic composition of coverage rather than holding it fixed.
The source-level Z-Score outcome and source clustering are unchanged. The
MSE-optimal bandwidth is re-selected for the total-effect column with topic
removed from the residualization, so the bandwidth remains optimal for the
specification actually estimated. The main specification here is Table 1 Model
5, which already uses all articles including military coverage, so both columns
sit on the same sample and the comparison is like-for-like. ‡ Cluster-robust SEs
from rdrobust.\\
\end{minipage}
\end{table}

